\section{Datasets}
% Describe MUTAG, ENZYMES, IMDB-MULTI: type of graphs, 
% nodes, labels.

% Any preprocessing steps.
The experiments in this work are conducted on three widely used benchmark datasets from the TUDataset collection: MUTAG, ENZYMES, and IMDB-MULTI.
These datasets represent diverse application domains and exhibit multiple graph sizes, densities and availability of node attributes, 
making them suitable for evaluating graph embedding methods.

\subsection{MUTAG}
MUTAG~\cite{debnath1991mutag} is a molecular graph dataset consisting of chemical compounds where each graph represents a molecule.Nodes correspond to atoms and edges represent chemical bonds between atoms. 
The task is to classify molecules into two classes according to their mutagenic effect on a bacterium.
The dataset contains a small number of graphs with discrete node labels representing atom types. 
Graphs in MUTAG are relatively small and sparse, making the dataset a common benchmark for evaluating graph classification methods on molecular data with labeled nodes. 

\subsection{ENZYMES}
The ENZYMES~\cite{rossi2015networkrepository} dataset is derived from protein tertiary structures and contains graphs representing enzymes. Nodes correspond to amino acids and edges indicate spatial proximity or interactions between residues. 
Each graph is labeled according to the enzyme’s functional class resulting in a multi class classification task.

\subsection{IMDB-MULTI}
MDB-MULTI~\cite{rossi2015networkrepository} is a social network dataset constructed from movie collaboration data. 
Each graph represents the ego network of actors appearing in a movie, nodes correspond to actors and edges indicate coappearance relationships. 
Unlike molecular datasets, IMDB-MULTI does not provide node labels or attributes.
